\section{Perpendicular standing spin wave and magnetic anisotropic study on amorphous FeTaC films - \href{https://doi.org/10.1109/TMAG.2016.2520981}{TMAG 7 52, 2520981 (2016)}}

\begin{description}
    \item[\underline{Perpendicular Standing spin-wave (PSSW):}]A perpendicular standing spin-wave is a magnetic excitation in a thin film where the spins oscillate in a standing-wave pattern across the film thickness, that is, perpendicular to the film surface.

    E.g: In a thick ferromagnetic film measured by FMR, higher-order resonance peaks appear because the magnetization forms standing waves between the two film surfaces.
\end{description}
\subsection*{Main points from the Paper}
\begin{enumerate}
    \item Magnetic anisotropy, spin wave (SW) excitation, and exchange stiffness constant of amorphous $\mathrm{Fe_{80}Ta_8C_{12}}$ $(d=20 - 200 \mathrm{nm})$ films were studied as a function of thickness using microstrip-FMR.
    \item The amorphous nature of the soft ferromagnetic (FM) alloys has no grain boundaries, which is important for the magnetization reversal and enhances the speed of spin-switching in spin transfer torque devices.
    \item In the case of FMR, the external microwave field couples with uniform and nonuniform SW modes. This excitation of magnetistatic modes occurs especially in dipole-dipole interaction dominated system where the wave vector $(\vec{q})$ lies in the plane.
    \item The modes with wave vector $(\vec{q})$ component perpendicular to the film surface is termed perpendicular standing spin-wave (PSSW) modes and the quantization of these modes is due to the geometrical confinement \cite{gui2007quantized}. 
    \item According to the Kittel model, the excitation of odd SW modes can be observed with the spin pinned by surface anisotropy interaction. However, in the case of partially surface pinned spins, both even and odd SW modes can be excited.
    \item The magnetic field sweep FMR spectra were carried out for different precessional frequencies and azimuthal angles.
    \item By increasing the thickness of the film $(\geq100 \mathrm{nm})$, $M-H$ curve clearly shows the transcritical loop manifesting the presence of stress-induced perpendicular anisotropy during the film deposition.
    \item The room temperature soft magnetic properties degrade drastically as the thickness increases due to the transition from the in-plane orientation of magnetization to the strip domain patterns.
    \item At the thickness of 20 nm a clear uniform precession mode without any other SWR mode in which the dynamic magnetization is uniform across thickness. 
    \item However, when the thickness is increased from 50 nm onward, the first PSSW mode is excited at lower absorption fields. The signal strength is 7 times smaller than the uniform mode.
    \item The absence of the higher order $(n>1)$ exchange-dominated PSSW modes could be due to very weak excitation, and the intensity of those nodes may be lower than our detection sensitivity.
    \item The increasing differences between the first PSSW modes and the uniform mode, could be understood on the basis of perpendicularly quantized SW vector, which is inversely proportional to the film thickness, i.e., $q=n\pi/d$.
    \item The angular dependence of resonance fields is governed by uniaxial magnetic anisotropy, in the case of 20 nm, $\phi_H$ dependence of $H_r$ for uniform mode at 6, 8 , and 10 GHz precessional frequencies.
    \item The typical field dependence of the resonant frequencies are shown in fig: 3(b). \hl{Watch this in detail}.
    \item The origin of in-plane uniaxial magnetic anisotropy in FeTaC thin films $(d=20$ and $50$ nm$)$ could be understood on the basis that the strong exchange coupling between the FM atoms plays an important role during the deposition process, which allows to form aligned FM atom pairs parallel to the film plane. 
    \item On further increase of the thickness, the in-plane uniaxial magnetic anisotropy disappears. 
\end{enumerate}
\subsection*{Summary}
The MS-FMR method was used to study magnetic anisotropy and spin-wave excitations in soft ferromagnetic FeTaC films. The angular variation of the resonance field within the film plane indicates a weak uniaxial anisotropy for films with thicknesses of 20 and 50 nm. Along with the uniform precession mode, the observation of the first perpendicular standing spin-wave mode confirms the role of both exchange and dipolar interactions. The exchange stiffness constant of these amorphous films was obtained from the in-plane angular behavior and from the frequency separation of the PSSW mode measured along the easy and hard magnetization directions. The extracted exchange stiffness increases from $2(4)\times10^{-7} \mathrm{to} 3.5(5)\times10^{-7}$ erg/cm as the film thickness increases from 50 to 200 nm.